\subsection{不可约表示的次数}

现在回到群 \(X(T)\) 的加法记号，并且不再假设 \(\rho\) 属于 \(X(T)\)。

定理 3. 最高权为 \(\lambda\) 的 G 的不可约表示空间的维数由下式给出

\[
\chi_ {\lambda} (e) = \prod_ {\alpha \in \mathrm{R} _ {+}} \frac {\langle \lambda + \rho , K _ {\alpha} \rangle}{\langle \rho , K _ {\alpha} \rangle}.
\]

置 \(\gamma = \frac{1}{2}\sum_{\alpha >0}K_{\alpha}\)，于是 \(\delta (\alpha)(\gamma) = 2\pi i\) 对每个单根 \(\alpha\) 成立（第 VI 章，第 1 节，第 10 号，命题 29）。直线 \(\mathbf{R}\gamma\) 不包含在任何超平面 \(\mathrm{Ker}\delta (\alpha)\) 中，所以 \(\exp (z\gamma)\) 对所有充分小的 \(z\in \mathbf{R}^*\) 都是正则元；对于所有 \(\mu \in \mathrm{X(T)}\) 和所有 \(z\in \mathbf{R}\)，有

\[
\mathrm{J} (\mu) (\exp (z \gamma)) = \sum_ {w \in \mathrm{W}} \varepsilon (w) e ^ {z \delta (\mu) (w ^ {- 1} \gamma)}.
\]

暂且接受下述引理：

引理 5. 有

\[
\mathrm{J} (\mu) (\exp (z \gamma)) = e ^ {z \delta (\mu) (\gamma)} \prod_ {\alpha > 0} (1 - e ^ {- z \delta (\mu) (K _ {\alpha})}).
\]

因此，函数 \(\mathrm{J}(\mu)(\exp (z\gamma))\) 是一个当 \(z\) 趋于 0 时趋于 1 的函数与下式的乘积

\[
z ^ {\mathrm{N}} \prod_ {\alpha > 0} \delta (\mu) (K _ {\alpha}) = (2 \pi i z) ^ {\mathrm{N}} \prod_ {\alpha > 0} \langle \mu , K _ {\alpha} \rangle
\]

其中 N = R+ 的基数。

首先假设 \(\rho \in \mathrm{X(T)}\)；应用定理 2 的推论 2，可知当 \(z\) 趋于 0 时，\(\chi_{\lambda}(z\gamma)\) 趋于

\[
\prod_ {\alpha > 0} \langle \lambda + \rho , K _ {\alpha} \rangle / \prod_ {\alpha > 0} \langle \rho , K _ {\alpha} \rangle ,
\]

因此在此情形下定理成立。

在一般情形下，只需注意，在证明定理 3 时总可以将 G 换成适当的连通覆盖，从而归结为前面的情形。

现在证明引理 5。令 \(z \in \mathbf{C}\)；将 \(\varphi_z\) 记为从 \(\mathfrak{t}\) 到 \(\mathbf{C}\)-代数 \(\operatorname{Map}(\mathrm{X}(\mathrm{T}), \mathbf{C})\) 的映射；这里的代数由从 \(\mathrm{X}(\mathrm{T})\) 到 \(\mathbf{C}\) 的映射组成；它将 \(H \in \mathfrak{t}\) 对应到映射

\[
\varphi_ {z} (H): \mu \mapsto \mu (\exp z H) = e ^ {z \delta (\mu) (H)}.
\]

有 \(\varphi_z(H + H') = \varphi_z(H)\varphi_z(H')\)，所以存在环同态

\[
\psi_ {z}: \mathbf {Z} [ \mathfrak {t} ] \to \operatorname{Map} (\mathrm{X} (\mathrm{T}), \mathbf {C})
\]

使得 \(\psi_z(e^H)(\mu) = e^{z\delta (\mu)(H)}\)。另一方面，根据第 VI 章第 3 节第 3 号命题 2，在 \(\mathbf{Z}[t]\) 中有如下关系：

\[
\sum_ {w \in \mathrm{W}} \varepsilon (w) e ^ {w \gamma} = e ^ {\gamma} \prod_ {\alpha > 0} (1 - e ^ {- K _ {\alpha}}).
\]

应用同态 \(\psi_z\)，并利用等式

\[
\psi_ {z} (e ^ {w \gamma}) (\mu) = e ^ {z \delta (\mu) (w \gamma)} = e ^ {z \delta (w ^ {- 1} \mu) (\gamma)},
\]

得到所述公式。

推论 1. 令 \(\| \|\) 为 \(\mathrm{X(T)}\otimes \mathbf{R}\) 上的范数。对于所有 \(\lambda \in \mathrm{X}_{++}\)，令 \(d(\lambda)\) 为 \(\mathrm{G}\) 的最高权为 \(\lambda\) 的不可约表示空间的维数。

\(a) \sup_{\lambda \in X_{++}} d(\lambda) / \| \lambda + \rho \|^{N} < \infty, \text{ 其中 } N = 1/2 (\dim G - \dim T).\)

\begin{enumerate}
\def\labelenumi{\alph{enumi})}
\setcounter{enumi}{1}
\item
  若 G 是半单的，则 \(\inf_{\lambda\in X_{++}}d(\lambda)/\|\lambda+\rho\|>0\)。
\item
  对于所有 \(\alpha \in \mathrm{R}_{+}\)，存在 \(\mathrm{A}_{\alpha} > 0\)，使得 \(|\langle \lambda + \rho, K_{\alpha} \rangle| \leq \mathrm{A}_{\alpha} \| \lambda + \rho \|\)，因而 \(d(\lambda) / \| \lambda + \rho \|^{\mathrm{N}} \leq \prod_{\alpha > 0} \mathrm{A}_{\alpha} / \langle \rho, K_{\alpha} \rangle\)。
\item
  假设 G 是半单的，将单根记为 \(\beta_{1},\ldots,\beta_{r}\)，并置 \(N_{i}=K_{\beta_{i}}\)。则
\end{enumerate}

\[
d (\lambda) \geq \prod_ {i = 1} ^ {r} \frac {\langle \lambda + \rho , N _ {i} \rangle}{\langle \rho , N _ {i} \rangle} = \prod_ {i = 1} ^ {r} \langle \lambda + \rho , N _ {i} \rangle ;
\]

由于 \(\langle \lambda +\rho ,N_i\rangle \geq \langle \rho ,N_i\rangle = 1\)，这意味着

\[
d (\lambda) \geq \sup _ {i} | \langle \lambda + \rho , N _ {i} \rangle |.
\]

若 G 是半单的，则 \(x \mapsto \sup |\langle x, N_{i} \rangle|\) 是 \(\mathrm{X(T)} \otimes \mathbf{R}\) 上的范数，必然与给定范数等价，故 b) 成立。

推论 2. 假设 G 是半单的；令 d 为整数。维数 \(\leq d\) 的 G 的表示的等价类集合是有限的。

推论 1 b) 蕴含，集合 \(X_{d}\) 中的元素 \(\lambda\) 属于 \(X_{++}\) 且满足 \(d(\lambda) \leq d\) 的元素所成的集合是有限的。对于所有 \(\lambda\) 属于 \(X_{d}\)，令 \(V_{\lambda}\) 为最高权为 \(\lambda\) 的不可约表示；每个维数 \(\leq d\) 的表示都同构于直和 \(\bigoplus_{\lambda \in X_{d}} V_{\lambda}^{n_{\lambda}}\)，其中 \(n_{\lambda} \leq d\)，故推论得证。

